\chapter{正文格式示例}

本章只演示标题、表格、插图、公式和注释的排版方式。
请在正式写作时删除所有占位说明，并用真实、可核验的论文内容替换。

\section{标题层级}

\subsection{三级标题}

\subsubsection{四级标题}

正文采用小四号宋体、首行缩进两个汉字和 1.5 倍行距。
默认的 \texttt{humanities} 类选项使用“一、”“（一）”“1.”“（1）”编号；
将主文件改为 \texttt{science} 选项后，标题改用阿拉伯数字分级编号。

\section{图表与公式}

表题置于表格上方。理工农医类默认按章编号，人文社科类默认全文连续编号。

\begin{table}[htbp]
  \centering
  \caption{表格格式示例}
  \label{tab:format-example}
  \begin{tabular}{lll}
    \toprule
    项目 & 内容 & 说明 \\
    \midrule
    示例一 & 占位内容 & 正式写作时替换 \\
    示例二 & 占位内容 & 不代表研究数据 \\
    \bottomrule
  \end{tabular}
\end{table}

\begin{figure}[htbp]
  \centering
  \fbox{\rule{0pt}{4cm}\rule{8cm}{0pt}}
  \caption{插图位置示例}
  \label{fig:format-example}
\end{figure}

公式居中排版，并按章连续编号。例如：
\begin{equation}
  y = f(x).
  \label{eq:format-example}
\end{equation}

此处演示文末集中注释标记\yibinnote{此条内容仅用于演示文末注释的排版位置。}。
